%! TEX root = ../main.tex
\documentclass[main]{subfiles}

\begin{document}
\section{実験}
    \label{実験}
先行実装及び先行研究に対する精度向上及び各要素の有効性を示すため，本研究では3つの実験を行う．\par 
(i) まず，先行実装（MT），先行研究（CMT），及び提案手法を対象とし，それぞれの識別性能を評価する．
(ii) 次に，MixStyleを既存モデルに直接適用する場合と，軽量CNNを介して統合する場合とで，精度を比較する．
(iii) 最後に，構成要素の比較を行う．半教師あり学習の三手法に対し，MixStyleを併用して学習を行い，上記の結果と合わせて各モジュールの導入効果を検証する．
なお，実験環境は以下の通り：
{\footnotesize
\begin{itemize}
    \item OS: Ubuntu 24.04.3 LTS (Noble Numbat) 
    \item Kernel: 6.14.0-36-generic
	\item CPU: AMD Ryzen 7 5700X 8-Core Processor, x86\_64
	\item GPU: NVIDIA GeForce RTX 4070, 12282 MiB
	\item RAM: 15GiB
\end{itemize}
}
\end{document}